If a user wishes to contribute a model to the \greensloth{} ecosystem, they can follow the steps outlined in this section. The process is designed to be straightforward and minimize effort for the user, while ensuring the model is well-documented and easily shareable with the community. The steps follow the same separation of power as \greensloth{}: \hyperref[sec:howto-creation]{1) Model Creation}, \hyperref[sec:howto-present]{2) Model Presentation}, and \hyperref[sec:howto-sharing]{3) Model Sharing}

\subsection{Model Creation}\label{sec:howto-creation}

The actual writing of the model cannot be assisted by \greensloth{}, as it is too complex a task. However, once a model is written in Python, using the \verb|mxlpy| package, \greensloth{}, can take over many tasks. The biggest is a bridge between the Python code and the actual documentation. To start, the user downloads the \greensloth[GreenSlothUtils] package (\url{https://github.com/ElouenCorvest/GreenSlothUtils}), and installs it as an editable package into their Python environment. This allows the user to use all the utilities of the \greensloth[GreenSlothUtils] package, including a step-by-step guide to creating a well-formatted model repository. The details of this process can be found in the \greensloth[GreenSlothUtils] documentation and in the materials and methods section~\cleverref{Section}{sec:greenslothutils} of this thesis. With the information and documentation complete, the user can then recreate the publication's figures in the already-generated template Jupyter Notebook. In the README file, these figures shall then be imported into the "Figures" section and shortly described. Due to rights and permissions, the user may not be able to share the original figure caption from the publication. The captions for these figures must describe not only what is visible, but also how the figures were obtained. The scientific context of the figures is not needed, as they are part of the publication. With the model now well packaged, validated, and documented, the user can proceed to the next step.

\subsection{Model Presentation}\label{sec:howto-present}

The model presentation is the recreation of the demonstrations later used to compare the model. As these demonstrations are standardized, the user does not need to worry about much code. They can use the already-generated Jupyter Notebook template and fill in the missing names of the quantities needed for each demonstration, all collected at the start of the notebook. If the current model is missing any required quantities, the user shall instead provide \verb|None|. In this case, each demonstration will decide what happens with the missing quantity. For example, if the model does not have a quantity for \gls{npq}, the demonstrations of the \gls{pam} protocol and the fitting will both calculate it from the \gls{F}. Sadly, in some cases, the demonstrations have specific rules or requirements. For example, the \gls{fvcb} add-on needs the \gls{vc} to be in a unit of \unit{\micro\mol\per\meter\squared\per\second}. If that is not the case in the current model, the user needs to either create a new quantity with the correct unit or let the demonstration handle it if the unit fits. These extra requirements are listed in the documentation of the \greensloth[GreenSlothUtils] package, where all the required functions are located. With the demonstrations now filled in, the user can add notes to the README file, if applicable. The README file will automatically import the figures generated by the demonstrations if the user followed the previous model-creation process. At this point, the model is well packaged and documented, and is now prepared to be shared with the community.

\subsection{Model Sharing}\label{sec:howto-sharing}

Once the model has reached the required state of being well packaged, validated, demonstrated, and documented, it can now be uploaded to the site. Before that, though, the user has to check whether their model overlaps with models already on the site. To do that, the main \greensloth{} repository contains two main glossaries in the overarching models directory. These two files contain all the rates and variables used in the models already in the ecosystem. The user has to go through these lists and see if any match their model. This is a tedious process; however, if there is a match, the user only has to include the respective Glossary ID in the correct column. Actually doing this during model creation is recommended. However, it had to be mentioned as the final obligation before the model could be shared. Once the user has done this, they can then create a pull request of their branch with the new model. The pull request will then be reviewed by the maintainers of \greensloth{} and accepted if all of the previous steps have been followed. If the pull request is accepted, the user does not need to do anything else, as the maintainers will handle the rest.

\subsection{After Sharing}

Since the user-uploaded model follows the required format, it can now be added to the website and made available to the community for use. To keep maintainers' workload as low as possible, this step has been automated as much as possible. Once the model directory is in the main \greensloth{} repository, the maintainers switch to the website branch (\url{https://github.com/ElouenCorvest/GreenSloth/tree/website}), and add the new model name inside a small \gls{json} file. This file predicts which models are extracted from the main repository and also serves as the storage for the tags. These tags also need to be added manually by the maintainer. Once the \gls{json} file is updated, the maintainer only needs to duplicate one of the other model \gls{html} files, rename it to the new model name, and that is it. The maintainer should then check on the website's dev environment that the new model is displayed correctly. If everything is done correctly, the rest of the work is done automatically. The maintainer can then build and deploy the website, making the new model available on the live website.